% s08-fisoc.tex —— §6 F-Isocrystals, p-divisible groups and modifications of vector bundles，原文 p.31–p.34
\srcpage{31}

\section{F-等晶体、\texorpdfstring{$p$}{p}-可除群与向量丛的修正}
\label{sec:fisoc}

本节给出一些补充，供 §7 与 §8 使用。

\subsection{\texorpdfstring{$Y$}{Y} 的紧化}
\label{sec:fisoc-compactification}

\subsubsection*{6.1.1.\ affinoid 情形}

设 $S=\Spa(R,R^+)$ 是 affinoid 完满空间。
回忆 §1 中的空间 $Y=Y_S$。
先设 $E=\FF_q((\pi))$。
此时 $Y=\mathbb{D}_S^*$。
这自然地延拓为 $\oo_E$-adic 空间
\[
\mathbb{D}_S\longrightarrow\mathbb{D}_{\FF_q}=\Spa(\oo_E),
\]
它无非是 $S$ 上的开圆盘，并且
\[
\mathbb{D}_S^*=\mathbb{D}_S\setminus\{\pi=0\}.
\]
但事实上，利用 $R$ 上由 $R^+$ 给出的整结构，这可以延拓为更大的 adic 空间
\[
\begin{aligned}
\mathcal{Y}&=\Spa\bigl(R^\circ\llbracket\pi\rrbracket,\;R^+ + R^\circ\llbracket\pi\rrbracket\bigr)_a\\
&=\Spa\bigl(R^\circ\llbracket\pi\rrbracket,\;R^+ + R^\circ\llbracket\pi\rrbracket\bigr)\setminus\{\pi=0\}\\
&\qquad\cup\Spa\bigl(R^\circ\llbracket\pi\rrbracket,\;R^+ + R^\circ\llbracket\pi\rrbracket\bigr)\setminus\{\varpi_R=0\},
\end{aligned}
\]
其中下标 “a” 表示取解析点，且 $\varpi_R\in R^{\circ\circ}\cap R^\times$。
回忆在此定义中，$R^\circ\llbracket\pi\rrbracket$ 带有 $(\varpi_R,\pi)$-adic 拓扑。
§1.1.2 的公式 (1) 给出的函数 $\delta$ 延拓为函数
\[
\delta:|\mathcal{Y}|\longrightarrow[0,1].
\]
空间 $\mathcal{Y}$ 是 $Y$ 在 $\oo_E$ 上的某种紧化，它通过添加除子 $(\pi)$ 与 $(\varpi_R)$ 得到，并且
\[
Y=\delta^{-1}\bigl({]0,1[}\bigr).
\]
记 $E^{1/p^\infty}$ 为 $E$ 的完满化的 $\pi$-adic 完备化。
有
\[
\mathcal{Y}\widehat{\otimes}_{\oo_E}\oo_E^{1/p^\infty}
=\Spa\bigl(R^\circ\llbracket\pi^{1/p^\infty}\rrbracket,\;R^+ + R^\circ\llbracket\pi^{1/p^\infty}\rrbracket\bigr)_a.
\]
这是 $\oo_E^{1/p^\infty}$ 上的完满空间。
开子集 $\{\pi\neq0\}$ 是完满域 $E^{1/p^\infty}$ 上的完满空间；$\{\varpi_R\neq0\}$ 经由 $T\mapsto\varpi_R$ 是完满域 $\FF_q((T^{1/p^\infty}))$ 上的完满空间。
在这个意义下，$\mathcal{Y}$ 是预完满的\yizhuw{预完满（preperfectoid）指：基变换到某个完满域扩张之后成为完满空间，而自身未必完满。}。

现设 $E|\QQ_p$。
事实是，同一个公式
\[
\begin{aligned}
\mathcal{Y}&=\Spa\bigl(\mathbf{A},[R^+]+\pi\mathbf{A}\bigr)_a\\
&=\Spa\bigl(\mathbf{A},[R^+]+\pi\mathbf{A}\bigr)\setminus V(\pi)
\;\cup\;\Spa\bigl(\mathbf{A},[R^+]+\pi\mathbf{A}\bigr)\setminus V([\varpi_R])
\end{aligned}
\]
\srcpage{32}
仍然定义一个 $\oo_E$-adic 空间。
这里 $\mathbf{A}=W_{\oo_E}(R^\circ)$，并且
\[
\mathbf{A}^+=\Bigl\{\textstyle\sum_{n\ge0}[x_n]\pi^n\in\mathbf{A}\ \Big|\ x_0\in R^+\Bigr\}.
\]
与前面一样，这是用除子 $(\pi)$ 与 $([\varpi_R])$ 对 $Y$ 作的紧化。
此外，若 $E_\infty$ 是由某个 Lubin--Tate 群的全部挠点生成的扩域的完备化，则 $\mathcal{Y}\widehat{\otimes}_{\oo_E}\oo_{E_\infty}$ 是完满的，其倾斜便是前面那个联系于 $\oo_{E_\infty}^\diamond$ 的等特征空间。

\subsubsection*{6.1.2.\ 一般情形}

回到等特征情形。
对区间 $I\subset[0,1]$（不同于 $\{0\}$ 与 $\{1\}$），记 $\mathcal{Y}_I$ 为 $\mathcal{Y}$ 中由半径函数 $\delta$ 定义的相应圆环。
公式
\[
\mathbb{D}_S=\mathcal{Y}_{S,[0,1[}\subset\mathcal{Y}_{S,]0,1[}=\mathbb{D}_S^*
\]
告诉我们：前述空间 $\mathcal{Y}_{S,]0,1[}$ 对任意 $\FF_q$-完满空间 $S$ 都可整体化，不必是 affinoid 完满的。
然而对一般的 $S$，$\mathcal{Y}_S$ 并非如此。
为此需要对 $S$ 取整模型。

更确切地说，考虑一个完满的 $\FF_q$-形式概形 $\mathcal{S}$，它有一个由仿射子集 $\mathrm{Spf}(A)$ 组成的开覆盖，其中 $A$ 对某个正则元 $\varpi_A\in A$ 是 $\varpi_A$-adic 的。
于是可以定义其一般纤维 $\mathcal{S}_\eta$ 为 $\FF_q$-完满空间。
对如上之 $A$，有 $\mathrm{Spf}(A)_\eta=\Spa\bigl(A[\tfrac{1}{\varpi_A}],A^+\bigr)$，其中 $A^+$ 是 $A$ 在 $A[\tfrac{1}{\varpi_A}]$ 中的整闭包（这几乎等于 $A$）。
现在给定这样的 $A$，便可定义一个 $\oo_E$-adic 空间
\[
\mathcal{Y}_{\mathcal{S}},
\]
使得
\[
Y_{\mathcal{S}_\eta}=\mathcal{Y}_{\mathcal{S},]0,1[}.
\]
这对等特征与不等特征的 $E$ 都成立。
我们不想过多深入细节，以下限于 affinoid 情形以简化。

\subsection{F-等晶体与向量丛}
\label{sec:fisoc-isocrystals}

现设 $E=\QQ_p$，并令 $S=\Spa(R,R^+)$ 是 affinoid 完满空间。
固定某个伪一致化元 $\varpi_R$。
晶体 $p$-adic Hodge 理论与曲线之间的联系由以下给出。
记 $\mathcal{Y}:=\mathcal{Y}_S$。
考虑环
\[
\mathrm{B}^+_{cris}(R^\circ/\varpi_R)
=H^0_{cris}\bigl(\Spec(R^\circ/\varpi_R)/\Spec(\ZZ_p),\mathcal{O}\bigr)\Bigl[\tfrac1p\Bigr],
\]
其中 $\mathcal{O}$ 是晶体位点上的结构层。
由于 $R^\circ/\varpi_R$ 上的 Frobenius 是满的，该位点有一个终对象，由 Fontaine 的环 $A_{cris}(R^\circ/\varpi_R)$ 给出。
这是
\[
W(R^\circ)\Bigl[\tfrac{[\varpi_R^n]}{n!}\Bigr]_{n\ge1}
\]
的 $p$-adic 完备化。
这是 Fontaine 的周期环：在 Witt 向量环上加入 $[\varpi_R^n]/n!$ 这类 divided power 元，使得 Frobenius 的作用可积\yizhuw{即 $A_{cris}$；加 divided power 结构是为了让 $p$-adic 对数与 Frobenius 在同一条链上相容，这正是晶体上同调能算出 $\mathrm{B}^+_{cris}$ 的原因。}。
于是有
\[
\mathrm{B}^+_{cris}=A_{cris}\bigl[\tfrac1p\bigr].
\]
现在的事实是，一个简单的计算给出：
\[
\Gamma\bigl(\mathcal{Y}_{[1/p^p,1]},\mathcal{O}\bigr)
\subset\mathrm{B}^+_{cris}(R^\circ/\varpi_R)
\subset\Gamma\bigl(\mathcal{Y}_{[1/p^{p-1},1]},\mathcal{O}\bigr).
\]
回忆 Frobenius 在 $\mathcal{Y}$ 上的作用满足
\[
\delta(\varphi(y))=\delta(y)^{1/p}.
\]
经晶体位点上的整体截面函子，$\Spec(R^\circ/\varpi_R)$ 上的 $F$-等晶体等同于一个有限型的射影模，它定义在 $\mathrm{B}^+_{cris}(R^\circ/\varpi_R)$ 上，并带有一个半线性自同态，其线性化是同构。
由此推出以下结论。

\begin{proposition}
$(\Spec(R^\circ/\varpi_R)/\Spec(\ZZ_p))_{cris}$ 上的 $F$-等晶体范畴，等价于 $\mathcal{Y}_{(R,R^+),]0,1[}$ 上 $\varphi$-等变向量丛的范畴。
\end{proposition}

\begin{remark}
最后我们看到，这个 $F$-等晶体范畴不依赖于 $\varpi_R$ 的选取。
但事实上，这一点此前就能看出，因为 $F$-等晶体的范畴在无穷小加厚下不变。
\end{remark}

\srcpage{33}
我们由此推出一个函子。

\begin{corollary}
从 $\mathcal{Y}_{]0,1]}$ 到 $Y=\mathcal{Y}_{]0,1]}\setminus V([\varpi_R])$ 的限制，诱导出一个函子
\[
\text{$F$-等晶体 on }\Spec(R^\circ/\varpi_R)\longrightarrow\Bun_{X_{R,R^+}}.
\]
\end{corollary}

这个函子总是全忠实的，但一般不是本质满的。
尽管如此，我们有下面的结果（[10]），后面不会用到它（尽管 Scholze 与 Weinstein 在 [30] 中用过，而我们将会用到 [30]）。

\begin{theorem}
当 $S$ 是一个几何点，也就是说 $S=\Spa(F)$ 且 $F$ 代数闭时，这是一个等价
\[
\text{$F$-等晶体 on }\Spec(\oo_F/\varpi_F)\xrightarrow{\ \sim\ }\Bun_{X_F}.
\]
\end{theorem}

当 $F$ 是完满的非代数闭域时，这个定理已经不成立。

\begin{remark}
必须极其小心定理 6.4 中的等价：它的逆不是正合函子，因此这并不是正合范畴的等价。
\end{remark}

\begin{hope}
曲线上的向量丛叠 $\Bun$ 是 $F$-等晶体的预叠关于 $\mathrm{Perf}_{\FF_p}$ 上忠实拓扑的叠化，也就是说，是预叠
\[
(R,R^+)\longmapsto \text{$F$-等晶体 on }\Spec(R^\circ/\varpi_R)
\]
所对应的叠。
这在某种意义上将说明，$\Bun$ 是 $F$-等晶体之叠的“一般纤维”，即刚性解析几何意义下的一般纤维。
\end{hope}

\subsection{\texorpdfstring{$p$}{p}-可除群与向量丛的修正}
\label{sec:fisoc-modifications}

设 $S=\Spa(R,R^+)$ 如前是 affinoid 完满空间。
现设 $H$ 是 $R^\circ/\varpi_R$ 上的 $p$-可除群。
考虑其协变 Dieudonné 晶体 $\mathbb{D}(H)_{\mathbb{Q}}$。
根据命题 6.1，可以联系于它一个向量丛
\[
\mathcal{E}\bigl(\mathbb{D}(H)_{\mathbb{Q}}\bigr),
\]
定义在 $X:=X_S$ 上。
现设 $S^\sharp=\Spa(R^\sharp,R^{\sharp,+})$ 是 $R$ 的一个 untilt。
存在同构
\[
R^{\sharp,\circ}/\varpi_R^\sharp\xrightarrow{\ \sim\ }R^\circ/\varpi_R,
\]
其中须把 $\varpi_R$ 换成 $\varpi_R^{1/p^n}$，$n\gg0$（回忆改变 $\varpi_R$ 不改变 $F$-等晶体 $\mathbb{D}(H)_{\mathbb{Q}}$，见注 6.2，故这一改变是无害的）。
记
\[
i:S^\sharp\hookrightarrow X_S
\]
为 $X_S$ 上相应的 Cartier 除子。
现设 $\widetilde{H}$ 是 $H$ 到 $R^{\sharp,\circ}$ 的提升。
此时特化
\[
i^*\mathcal{E}\bigl(\mathbb{D}(H)_{\mathbb{Q}}\bigr)
\]
与 $\widetilde{H}$ 的万有扩张的李代数等同，因此带有 Hodge 滤过
\[
0\longrightarrow\omega_{\widetilde{H}_D}\bigl[\tfrac1p\bigr]
\longrightarrow i^*\mathcal{E}\bigl(\mathbb{D}(H)_{\mathbb{Q}}\bigr)
\longrightarrow\Lie(\widetilde{H})\bigl[\tfrac1p\bigr]\longrightarrow0.
\]
下面是 Fontaine 关于 $p$-可除群的比较定理的内容。

\begin{proposition}
设 $V_p(\widetilde{H})$ 是 $S$ 上与 $\widetilde{H}$ 的一般纤维相联系的 pro-étale $\QQ_p$-局部系统，而 $V_p(\widetilde{H})\otimes_{\QQ_p}\mathcal{O}_{X_S}$ 是 $X_S$ 上相应的斜率 $0$ 半稳定向量丛。
此时存在正合列
\[
0\longrightarrow V_p(\widetilde{H})\otimes_{\QQ_p}\mathcal{O}_{X_S}
\longrightarrow\mathcal{E}\bigl(\mathbb{D}(H)_{\mathbb{Q}}\bigr)(1)
\longrightarrow i_*\omega_{\widetilde{H}_D}\bigl[\tfrac1p\bigr]\longrightarrow0.
\]
\end{proposition}

这里用 $\mathcal{O}_{X_S}(1)$ 作扭转，原因在于 $\mathbb{D}(\widetilde{H}_{\mathbb{Q}})$ 是协变 Dieudonné 模\yizhuw{原文此句作“the fact that we $\mathbb D(\widetilde H_{\mathbb Q})$ is the covariant Dieudonné-module”，多出一个“we”，显系笔误；译文按通顺读法处理。}。
现在我们可以把它翻译成 Hecke 叠的语言了。

\begin{definition}
我们把 $\Spa(\QQ_p)^\diamond$ 上的 $\mathcal{BT}^{d,h}$ 定义为 $\QQ_p$-affinoid 完满空间上的预叠\yizhuw{预叠（prestack）指取值于群胚、但不要求下降的函子；$\mathcal{BT}$ 一般不满足下降，故只能算预叠，后面要靠它过渡到真正的叠。}。对 $\QQ_p$-affinoid 完满空间 $(R,R^+)$，$\mathcal{BT}(R,R^+)$ 定义为 $R^\circ$ 上高度为 $h$、维数为 $d$ 的 $p$-可除群的群胚。
\end{definition}

当然，$\mathcal{BT}^{d,h}$ 不是叠。
我们将把它当作中间预叠，用来构造两个叠之间的态射。

\begin{remark}
我们本可以取 $R^+$ 上的 $p$-可除群，而不是 $R^\circ$ 上的。
正如前面（3.2）所说，我们关心的对象都是 partially proper 的，因此要求 $\mathcal{BT}$ 是 partially proper 的并无害处\yizhuw{partially proper 指 $\Spa$ 意义下的“部分紧”（proper 性的开覆盖版本），此处保留英文。}。
\end{remark}

\srcpage{34}
我们由此推出以下结论\yizhuw{原文作“We thus deduces the following”，动词形式有误，译文按“由此推出”处理。}。

\begin{corollary}
取 $G=\mathrm{GL}_h$，$\mu(z)=\mathrm{diag}(\underbrace{z,\dots,z}_{d\text{-times}},1,\dots,1)$。
存在态射
\[
u:\mathcal{BT}^{d,h}\longrightarrow Hecke^\mu,
\]
它由与某个 $p$-可除群的 $F$-等晶体相联系的向量丛的 Hodge 滤过给出。
它满足：
\begin{enumerate}
\item[(1)] $\overleftarrow{h}\circ u$ 由与该 $p$-可除群模 $\varpi_R$ 的归约的 $F$-等晶体相联系的、经 $\mathcal{O}(1)$ 扭转后的向量丛给出；
\item[(2)] $\overrightarrow{h}\circ u$ 由与该 $p$-可除群的一般纤维相联系的 $\QQ_p$-pro-étale 局部系统相联系的斜率 $0$-半稳定向量丛给出（即其有理 Tate 模）。
\end{enumerate}
\end{corollary}

还有另一种“对偶”的方式，用 Hodge--Tate 周期来描述这个修正。
事实上，沿用前面的记号，存在 pro-étale 层的 $\QQ_p$-线性态射
\[
\alpha_{\widetilde{H}}:V_p(\widetilde{H})\longrightarrow\omega_{\widetilde{H}_D}\bigl[\tfrac1p\bigr].
\]
关于这个态射在赋值环上的定义与性质，例如可参看 [8] 的 §5。
回忆以下事实（其证明与 [8] 的定理 2（§5.3.2）相同，因为 Faltings 的整比较定理在任何完满环上都成立，见 [3]）。

\begin{proposition}
在 $S^\sharp$ 上存在向量丛的正合列
\[
0\longrightarrow\omega^\vee_{\widetilde{H}}\bigl[\tfrac1p\bigr]\otimes\QQ_p(1)
\xrightarrow{\ \alpha^\vee_{\widetilde{H}_D}(1)\ }
V_p(\widetilde{H})\otimes_{\QQ_p}\mathcal{O}_{S^\sharp}
\xrightarrow{\ \alpha_{\widetilde{H}}\ }
\omega_{\widetilde{H}_D}\bigl[\tfrac1p\bigr]\longrightarrow0.
\]
\end{proposition}

在前一个命题中，
\[
V_p(\widetilde{H})\otimes\mathcal{O}_{S^\sharp}
=i^*\bigl(V_p(\widetilde{H})\otimes\mathcal{O}_{X_S}\bigr),
\]
这是 $S^\sharp$ 上的向量丛。

\begin{proposition}
态射 $u:\mathcal{BT}^{d,h}\to Hecke^\mu$ 由 $V_p(\widetilde{H})\otimes_{\QQ_p}\mathcal{O}_{X_S}$ 的修正诱导：该修正由 $i^*\bigl(V_p(\widetilde{H})\otimes_{\QQ_p}\mathcal{O}_{X_S}\bigr)$ 的 Hodge--Tate 滤过给出，而后者联系于命题 6.11 的 Hodge--Tate 正合列。
\end{proposition}

于是，若着眼于 $\overleftarrow{h}$ 或 $\overrightarrow{h}$，我们对态射 $u$ 便有两种等价的看法：
\begin{itemize}
\item 若着眼于 $\overleftarrow{h}$，它由 de Rham 周期定义；
\item 若着眼于 $\overrightarrow{h}$，它由 Hodge--Tate 周期定义。
\end{itemize}
